The experiments address four questions: whether the stored cuts are valid in
practice, and what certification prevents; whether the cuts change SCIP's
results on benchmark models, including models that SCIP finds hard; what the
separator does and what it costs; and how much of the root gap the cuts close
on instances with the structure of Sections~\ref{sec:composition}
and~\ref{sec:quadratic}, compared with SCIP's own alternatives and with Gurobi.
We report five campaigns. Each was run with solver code, limits and protocol
fixed before its first run; we call these results \emph{prospective}, and we
label analyses designed after results were seen as \emph{post hoc}.
Campaigns 1 and~2 used earlier implementations and are summarized in
Appendix~\ref{app:campaigns12}. Campaign~3 tests the final implementation on
MINLPLib models (Parts~3A and~3B) and on a constructed family (Part~3C); a
post hoc diagnostic of its results is Part~3P. Campaign~4 was designed after
campaign~3 to answer the questions that its results raised: it adds
comparators, fresh instances, larger models, a variant of the family in which
the block minima do not give the optimum, and an ablation of the certificate;
the job builder and the reference computation of Part~4C4 were written during
campaign~4, as its protocol specifies, before any run of that part.
Campaign~5 was designed after campaign~4: it runs the star algorithm in the
solver, adds the remaining comparators on the fresh instances, tests larger cut
limits and compares the certificate with a numerical global solve.
Table~\ref{tab:overview} lists the parts.

\subsection{Setup}\label{sec:setup}

All runs use SCIP~10.0.2 through PySCIPOpt~6.2.1 and Python~3.13, with one
SCIP thread and one BLAS thread per process, on a Linux host (WSL2) with 18
cores and 36 hardware threads that was shared with unrelated jobs. The
one-minute load average was 14 to 21 during the prospective runs of
campaign~3 (2 to 21 during Part~3P), 2 to 9 during campaign~4 (5th to 95th
percentile; at most 15, in Part~4C4) and 1 to 20 during campaign~5, apart from
a spike to 126 during Part~5S (Section~\ref{sec:star-bench}). Within a
part, at most six runs were active at a time, and the modes of one instance
ran back to back in rotated order; the Part~3P runs on the path family
overlapped the last 28 minutes of Part~3A, when up to ten runs were active. We
compare times only within these paired runs of one part. Part~4C2 has no
baseline of its own and is compared with the campaign-3 runs of the same
instances; root bounds of identical runs repeated exactly across campaigns
(all 100 baseline root bounds rerun in Part~3P), whereas solved counts and
times are not paired across campaigns, and we say so where it matters. Timings
are descriptive. Times are in-process wall times charged to the time budget
(model build, discovery, separation and solve), excluding interpreter start-up
and the primal check; \emph{SCIP time excluding the callback} is SCIP's solving
time minus the time spent in the separator callback. Full runs have a budget of
300\,s; root runs have a node limit of one and a budget of 60\,s (Parts~3A, 3B,
4B2) or 120\,s (path family, Parts~4D and~5S). In Parts~3A and~3B a seed is
SCIP's random seed shift; constructed instances are identified by their
generator seed and run with SCIP's seed~0.

We write \emph{SCIP} for SCIP with default settings, \emph{SCIP-nolocks} for SCIP
with its implicit-discreteness presolve step switched off
(\code{constraints/nonlinear/checkvarlocks = d}), and \emph{SCIP-extra} for SCIP
with its disabled-by-default nonconvex separators switched on
(Section~\ref{sec:related}); all three use our model builder and add no
certified cut. Gurobi~13.0.3 \citep{Gurobi13} solves the same model with one
thread, its nonconvex mode (\code{NonConvex=2}) and the same gap limit. The
separator modes are listed in Table~\ref{tab:modes}. Models come from MINLPLib
\citep{BussieckDrudMeeraus2003,Vigerske2026MINLPLib} in OSiL format, the XML
instance format of Optimization Services (Appendix~\ref{app:instances}),
stratified by the library's metadata into convex models, nonconvex models
without integer variables, and nonconvex models with integer variables.

A model counts as \emph{solved} if SCIP (or Gurobi) reports optimality or the
relative gap limit $10^{-4}$, and the returned incumbent passes an independent
check of all original rows, bounds, domains, integrality restrictions and the
objective at a scaled tolerance of $10^{-5}$; the checker evaluates the source
expressions without the model builder. Times are compared by shifted geometric
means (shift 1\,s) and by the median of per-run ratios, over runs solved in
every compared mode. Dual bounds $z$ and $z'$ are compared with the tolerance
$10^{-4}\max\{1,\abs z,\abs{z'}\}$, the gap limit. The \emph{root gap closed} by
a mode is $(z_{\text{mode}}-z_{\text{base}})/(z^\ast-z_{\text{base}})$ for
minimization, where $z$ denotes root dual bounds and $z^\ast$ the optimal or
best known value; for a run that ends at the root without reporting a root
bound, because the root node was pruned, $z$ is its final dual bound. Every
recorded cut is replayed in a fresh process (Section~\ref{sec:replay}). A run
whose worker failed has no complete cut log; its cut count is reported as
unknown, not as zero. The protocols also prespecify statistics that we do not
report here, such as means per seed, bound comparisons at the tolerance
$10^{-6}$ and scan counts per stratum; they are in the supplementary archive and
do not change the conclusions.

\begin{table}[t]
\centering\footnotesize
\caption{Parts of campaigns 3 to 5. Full runs have a budget of 300\,s; root
runs have a node limit of one. All parts except 3P are prospective. Modes are
defined in Appendix~\ref{app:software} (Table~\ref{tab:modes}).}\label{tab:overview}
\begin{tabular}{@{}l>{\raggedright\arraybackslash}p{0.32\textwidth}>{\raggedright\arraybackslash}p{0.43\textwidth}l@{}}
\toprule
Part & Instances & Modes & Runs\\
\midrule
3A & 30 MINLPLib models of campaign 2; full runs seeds 0--2, root runs seed 0 & SCIP, all, auto; root runs also all-diag & full, root\\
3B & 30 structure-selected models; full runs seeds 0--1 & SCIP, all, auto; root runs also all-diag & full, root\\
3C & 20 path-family instances & SCIP, all-diag-mech & full, root\\
3P & 3A, 3B (root runs); 3C (post hoc) & whole-row variant of the cut modes; 3C also wide limits & root; full, root\\
\midrule
4C2 & the 20 instances of 3C & SCIP-nolocks, SCIP-extra, frozen-wide, Gurobi & full, root\\
4C3 & 20 fresh path-family instances & SCIP, all-diag-mech, rowdir-wide, SCIP-extra, Gurobi & full, root\\
4C4 & the instances of 4C3 with a binding row & SCIP, frozen-wide, rowdir-wide, SCIP-extra, Gurobi & full, root\\
4B2 & the 30 models of 3B & no aggregation: SCIP, all, all-diag, all-diag-rowdir; SCIP-extra & root\\
4D & 20 larger models that SCIP finds hard & full: SCIP, all, auto, SCIP-extra; root: SCIP, all, all-diag, all-diag-rowdir, SCIP-extra & full, root\\
4U & recorded cuts of campaigns 3 and 4 & uncertified constants (samples, local solver) & offline\\
4S & random stars with $10$ to $10^4$ leaves & two star algorithms & offline\\
\midrule
5S & 30 star-family instances & SCIP, SCIP-nolocks, SCIP-extra, rowdir-star4, agg-star4, agg-star, Gurobi & full, root\\
5C-a & the instances of 4C3 and 4C4 & SCIP-nolocks & full, root\\
5C-b & the instances of 4C4 & larger cut caps (32n, 64n), both direction rules & root\\
5U & the cuts of 4U & numerical global solve (Gurobi) & offline\\
\bottomrule
\end{tabular}
\end{table}
